% Resposta da questão 2

Para comparar a densidade de massa e a viscosidade da manteiga de amendoim, mel e água, apresentamos valores típicos encontrados na literatura:

\begin{itemize}
	\item \textbf{Manteiga de amendoim:}
	\begin{itemize}
		\item Densidade: $\approx 1200$ kg/m$^3$
		\item Viscosidade: $\approx 250{,}000$ a $350{,}000$ mPa·s (muito alta, depende da temperatura e formulação)
	\end{itemize}
	\item \textbf{Mel:}
	\begin{itemize}
		\item Densidade: $\approx 1400$ kg/m$^3$
		\item Viscosidade: $\approx 2{,}000$ a $10{,}000$ mPa·s (varia com a temperatura)
	\end{itemize}
	\item \textbf{Água:}
	\begin{itemize}
		\item Densidade: $\approx 1000$ kg/m$^3$ (a 20$^\circ$C)
		\item Viscosidade: $\approx 1$ mPa·s (a 20$^\circ$C)
	\end{itemize}
\end{itemize}

\textbf{Comparação:}

A manteiga de amendoim apresenta a maior viscosidade entre as três substâncias, sendo extremamente espessa e resistente ao escoamento. O mel possui viscosidade intermediária, ainda muito superior à da água, mas bem menor que a da manteiga de amendoim. A água, por sua vez, tem viscosidade muito baixa, fluindo facilmente. Em relação à densidade, o mel é o mais denso, seguido pela manteiga de amendoim e, por último, a água.
